% Exact research-candidate source-aware proof for the toric fixed-point scope of P06.1.
\begin{proof}[Proof of \Cref{thm:p061-toric-fixed-point-constancy}]
We prove separately the affinoid algebra assertion, the neighbourhood
assertion, and the pluriharmonic assertion.  This also records why neither
simpliciality nor a characteristic hypothesis enters the argument.

\smallskip
\noindent
\emph{The positive integral weight.}
We first recall directly why the rational polyhedral semigroup
\(S=\sigma^{\vee}\cap M\) is finitely generated.  Choose lattice vectors
\(e_1,\ldots,e_q\) generating the rational cone \(\sigma^{\vee}\).
For \(u\in S\), write \(u=\sum_j b_je_j\) with \(b_j\geq0\), subtract
\(\sum_j\lfloor b_j\rfloor e_j\), and observe that the remainder belongs
to
\[
 M\cap\left\{\sum_j c_je_j:0\leq c_j\leq1\right\}.
\]
This is a finite set.  The \(e_j\) together with these finitely many
remainders therefore generate \(S\).  Choose a system of nonzero monoid
generators \(u_1,\ldots,u_r\).  Since \(\sigma\) is full-dimensional,
\(S\cap(-S)=\{0\}\): an element in this intersection vanishes on the
spanning cone \(\sigma\).  Consequently
\[
 \mathfrak m_{\sigma}:=
 \bigoplus_{u\in S\setminus\{0\}}K\chi^u
\]
is a monomial maximal ideal with quotient \(K\); it defines the stated
torus-fixed point \(x_{\sigma}\).  An integral interior vector exists:
for example, the sum of the primitive lattice generators of all rays of
\(\sigma\) lies in \(\operatorname{int}(\sigma)\).  Since
\(v\in\operatorname{int}(\sigma)\),
\begin{equation}
\label{eq:p061-toric-positive-weight}
       \ell(u)=\langle u,v\rangle>0
       \qquad\text{for every }u\in S\setminus\{0\}.
\end{equation}
Indeed, if a nonzero \(u\in\sigma^{\vee}\) vanished at the interior point
\(v\), then \(u\) would vanish on a neighbourhood of \(v\) inside
\(N_{\mathbf R}\), and hence \(u=0\).  Write
\(w_i=\ell(u_i)\in\mathbf Z_{>0}\).  If
\(u=\sum_i a_i u_i\), then multiplicativity of a seminorm shows that the
inequalities in \eqref{eq:p061-toric-weighted-affinoid} for the \(u_i\)
imply the one for \(u\).  Thus the finite and infinite descriptions of
\(U_n\) agree.

\smallskip
\noindent
\emph{The strict affinoid algebra.}
The monoid generators give a closed immersion
\[
 X_{\sigma}\longrightarrow\mathbf A^r_K,
 \qquad T_i\longmapsto\chi^{u_i}.
\]
The subset \(U_n\) is the inverse image of the closed polydisc
\[
          |T_i|\leq |\varpi|^{nw_i}\qquad(1\leq i\leq r).
\]
All its polyradii belong to \(|K^\times|\), so this is a strict affinoid
domain of \(\mathbf A_K^{r,\mathrm{an}}\).  Its inverse image under the
closed immersion is the corresponding closed analytic subspace, hence an
affinoid domain of \(X_{\sigma}^{\mathrm{an}}\).
After putting
\begin{equation}
\label{eq:p061-toric-rescaled-monomial}
        Z^u:=\varpi^{-n\ell(u)}\chi^u,
\end{equation}
its affinoid algebra is
\begin{equation}
\label{eq:p061-toric-semigroup-tate-algebra}
 \mathscr A_n\simeq K\langle S\rangle
 :=\left\{\sum_{u\in S}a_uZ^u:
       a_u\in K,\ \text{and }|a_u|\longrightarrow0
       \text{ off finite subsets of }S\right\}.
\end{equation}
Here the norm is
\begin{equation}
\label{eq:p061-toric-semigroup-gauss-norm}
       \left\|\sum_{u\in S}a_uZ^u\right\|_{\mathrm G}
                 :=\max_{u\in S}|a_u|.
\end{equation}
For completeness, let \(I_S\subset K[T_1,\ldots,T_r]\) be the toric
ideal.  It is generated by the binomials \(T^a-T^b\) for which
\(\sum_i a_i u_i=\sum_i b_i u_i\).  The two monomials in each such
binomial have the same \(\ell\)-weight.  Hence the substitution
\(T_i=\varpi^{nw_i}Z_i\) carries the corresponding relation to a nonzero
scalar multiple of \(Z^a-Z^b\).  It follows that the affinoid algebra of
the displayed polydisc section is
\[
 K\langle Z_1,\ldots,Z_r\rangle/I_SK\langle Z_1,\ldots,Z_r\rangle,
\]
which is the algebra in \eqref{eq:p061-toric-semigroup-tate-algebra}.
This also proves that it is strictly \(K\)-affinoid.  One may verify the
last identification directly: choose for every \(u\in S\) one expression
\(u=\sum_i a_i(u)u_i\).  By \eqref{eq:p061-toric-positive-weight}, a
set \(\{u\in S:\ell(u)\leq L\}\) is finite for every \(L\), because
\(\ell(u)=\sum_i a_i(u)w_i\) and all \(w_i\) are positive.  Hence a
sequence escaping all finite subsets of \(S\) has
\(\sum_i a_i(u)\to\infty\), so a convergent semigroup series lifts to a
Tate series in the \(Z_i\), with exactly the same Gauss norm.  Since the
quotient map is contractive, the semigroup Gauss norm is at most the
quotient norm; the chosen same-norm lift gives the reverse inequality.
Thus the quotient norm is exactly the norm in
\eqref{eq:p061-toric-semigroup-gauss-norm}.

We next determine the supremum norm and canonical reduction rather than
merely appealing to the presentation.  The norm
\eqref{eq:p061-toric-semigroup-gauss-norm} is multiplicative.  To see
this, normalize two nonzero series \(f,g\) to have Gauss norm one.  Their
coefficientwise reductions are nonzero elements of
\(\widetilde K[S]\).  The latter ring is a domain because \(S\subset M\)
and hence \(\widetilde K[S]\subset\widetilde K[M]\).  Thus the product of
the two reductions is nonzero, which gives
\(\|fg\|_{\mathrm G}=1\).  Scaling proves multiplicativity in general.
In particular, the Gauss norm itself defines a point
\(\eta_n\in\mathcal M(\mathscr A_n)\).  Every point of
\(\mathcal M(\mathscr A_n)\) is bounded by the spectral norm, and the
power trick together with multiplicativity gives, for every \(f\),
\[
       |f(x)|\leq\|f\|_{\mathrm G}=|f(\eta_n)|.
\]
Consequently the Gauss norm is the supremum norm on \(U_n\).

It follows coefficientwise that the canonical reduction is
\begin{equation}
\label{eq:p061-toric-canonical-reduction}
             \widetilde{\mathscr A_n}\simeq\widetilde K[S].
\end{equation}
Indeed, a Gauss-bounded series has coefficients in \(K^{\circ}\), only
finitely many of which survive modulo \(K^{\circ\circ}\), and its image
vanishes precisely when its Gauss norm is strictly less than one.  Since
\(\widetilde K[S]\) is a domain, its spectrum has exactly one generic
point.  Berkovich's description of the Shilov boundary by the generic
points of the canonical reduction
\cite[Proposition~2.4.4]{Berk12} therefore gives
\begin{equation}
\label{eq:p061-toric-singleton-shilov}
                    \Gamma(U_n)=\{\eta_n\}.
\end{equation}
There is no hidden uniqueness issue here: the preceding norm calculation
already shows directly that the singleton \(\{\eta_n\}\) is a boundary.
In the original toric coordinates it is characterized by
\begin{equation}
\label{eq:p061-toric-weighted-gauss-point}
       |\chi^u(\eta_n)|=|\varpi|^{n\ell(u)}
       \qquad(u\in S).
\end{equation}

\smallskip
\noindent
\emph{The neighbourhood basis.}
The fixed point \(x_{\sigma}\) belongs to the ordinary interior of every
\(U_n\), because
\(|\chi^{u_i}(x_{\sigma})|=0<|\varpi|^{nw_i}\) for all \(i\).  Moreover,
\(U_{n+1}\subset U_n\), and
\begin{equation}
\label{eq:p061-toric-intersection}
                    \bigcap_{n\geq1}U_n=\{x_{\sigma}\}.
\end{equation}
Indeed, a point in the intersection annihilates every \(\chi^{u_i}\),
and hence every nonconstant monomial in \(K[S]\); its seminorm is therefore
the evaluation seminorm at \(x_{\sigma}\).

We now prove the required cofinality, which is stronger than merely saying
that each \(U_n\) contains \(x_{\sigma}\).  Let \(V\) be any ordinary open
neighbourhood of \(x_{\sigma}\).  If no \(U_n\) were contained in \(V\),
then
\[
          F_n:=U_n\cap(X_{\sigma}^{\mathrm{an}}\setminus V)
\]
would be a nested sequence of nonempty closed subsets of the compact
affinoid space \(U_1\).  Their intersection would be nonempty.  This
contradicts \eqref{eq:p061-toric-intersection}, since
\(x_{\sigma}\in V\).  Hence \(U_n\subset V\) for all sufficiently large
\(n\).  Together with the ordinary interior observed above, this proves
that the \(U_n\) form a basis of strict affinoid neighbourhoods.

\smallskip
\noindent
\emph{Maximum and minimum.}
Choose \(n\) sufficiently large that \(U_n\subset\Omega\).  Classical
pluriharmonicity restricts along the analytic-domain inclusion
\(U_n\hookrightarrow\Omega\), either by sheafiness or by functoriality
\cite[Proposition~7.2(1),(7)]{GR25Classical}.  Hence
\(h|_{U_n}\) is pluriharmonic.  The affinoid maximum--minimum statement
\cite[Proposition~7.3]{GR25Classical}, obtained by applying the affinoid
maximum principle \cite[Theorem~6.4]{GR25Classical} to both \(h\) and
\(-h\), and \eqref{eq:p061-toric-singleton-shilov} give
\[
 \max_{U_n}h=h(\eta_n)=\min_{U_n}h.
\]
Thus \(h\) is constant on \(U_n\).  The same argument applies to every
larger \(n\), which proves the asserted local constancy.

Finally, local constancy already gives both advertised regularity
conclusions.  Put
\[
 W_n:=\{x\in X_{\sigma}^{\mathrm{an}}:
           |\chi^{u_i}(x)|<|\varpi|^{nw_i}\text{ for }1\leq i\leq r\}.
\]
This is an ordinary open neighbourhood of \(x_{\sigma}\) contained in
\(U_n\).  If \(h|_{U_n}=c\), then \(h|_{W_n}=c\).  Choose
\(a\in K^{\times}\) with
\(|a|\ne1\).  Then
\begin{equation}
\label{eq:p061-toric-constant-unit-log}
             c=\frac{c}{\log|a|}\log|a|,
\end{equation}
so \(h|_{U_n}\) is \(\mathbf R\)-PL by
\cite[Definition~5.1]{GJR26}.  Since both signs are classically psh,
\cite[Theorem~4.3]{GR25Classical} makes both signs semipositive; this is
harmonicity by \cite[Definition~7.2]{GJR26}.  Equivalently,
\eqref{eq:p061-toric-constant-unit-log} is a real multiple of the
logarithm of an analytic unit, which is harmonic by
\cite[Remark~7.3(5),(6)]{GJR26}.

Every algebraic step used only that \(S\) is a pointed affine submonoid of
a lattice and that the positive integral weight \(\ell\) exists.  No ray
independence, finite quotient, averaging, or division by an integer was
used.  Thus neither simpliciality nor any characteristic restriction is
present.
\end{proof}
